\section*{第 \ref{ch:b2:field-lab-methods-2} 章——野外与实验方法 II}
% numbers from the Book 4 second-writer scratch calculator calc/ch31.py

\begin{solution}{exo:b2:field-lab-methods-2:1}
(a) (i) \omterm{def:b2:field-lab-methods-2:textures}{辉绿结构}；(ii) \omterm{def:b2:field-lab-methods-2:textures}{粒状变晶结构}；(iii) \omterm{def:b2:field-lab-methods-2:porphyroblast}{包含变晶}，即充满包裹体的\omterm{def:b2:field-lab-methods-2:porphyroblast}{变斑晶}；
(iv) \omterm{def:b2:field-lab-methods-2:porphyroblast}{反应边}，即包围正在分解的石榴子石的边缘。(b) 由
\eqref{eq:b2:field-lab-methods-2:lami} 得 $\gamma_1/\gamma_2 = \sin\ang{104}/\sin\ang{128} =
0.9703/0.7880 = 1.23$：该界面的拉力比另外两个大 \qty{23}{\percent}。
\end{solution}

\begin{solution}{exo:b2:field-lab-methods-2:2}
(a) $Q + A + P = 80$：Q $= 100 \times 25/80 = 31.3$，A $= 18.8$，P $= 50.0$。(b) P/(A+P) $=
40/55 = \qty{72.7}{\percent}$，介于 65 与 90 之间，且 Q 介于 20 与 60 之间：花岗闪长岩。
(c) $M' = 12 + 8 = \qty{20}{\percent}$，低于 35：浅色岩。
\end{solution}

\begin{solution}{exo:b2:field-lab-methods-2:3}
(a) $\sigma = \sqrt{0.12 \times 0.88/800} = 0.0115$，即 $12.0 \pm 1.1$ 个百分点，相对
误差为 $0.0115/0.12 = \qty{9.6}{\percent}$。(b) 点数增至四倍，误差减半：$0.0057$。
\end{solution}

\begin{solution}{exo:b2:field-lab-methods-2:4}
(a) 失重为 \qty{0.0288}{g}，且 $0.0288/0.8000 = \qty{3.60}{\percent}$。
(b) $0.7712/8.7712 = 0.0879$。
(c) 灼烧后粉末：$5.20/0.08792 = \qty{59.1}{\percent}$；原始粉末：
$59.14 \times (1 - 0.0360) = \qty{57.0}{\percent}$。
\end{solution}

\begin{solution}{exo:b2:field-lab-methods-2:5}
(a) $\sigma = 0.10 \times 0.02 = 0.002$，$n = 0.02 \times 0.98/0.002^2 = 4900$ 个点。
(b) $\sigma = 0.03$，$n = 0.30 \times 0.70/0.03^2 = 233.3$，故取 234 个点。(c) 相对误差固定时，
$n = (1-p)/(p\,\varepsilon^2)$ 与 $1/p$ 成正比：副矿物所需点数为原来的 21 倍，
超过一般薄片所能提供的点数。
\end{solution}

\begin{solution}{exo:b2:field-lab-methods-2:6}
(a) RGM-1：$72.89/73.45 = \qty{99.2}{\percent}$；BCR-2：$54.17/54.1 = \qty{100.1}{\percent}$。
(b) RGM-1：$0.19/\sqrt{20} = 0.042$，且 $-0.56/0.042 = -13.2$；BCR-2：$0.13/\sqrt{20} = 0.029$，且
$0.07/0.029 = 2.4$。(c) RGM-1 存在系统偏差：其均值与推荐值相差
自身标准差的 13 倍，远超偶然误差所能解释的范围（若推荐值正确）。BCR-2 处于偶然误差允许的边界。
RGM-1 的精密度并不差：其相对标准差
$0.19/72.89 = \qty{0.26}{\percent}$，与 BCR-2 的 \qty{0.24}{\percent} 接近。问题在于
正确度，而非精密度。
\end{solution}

\begin{solution}{exo:b2:field-lab-methods-2:7}
(a) $M = 39.098 + 3 \times 24.305 + 26.982 + 3 \times 28.085 + 12 \times 15.999 + 2 \times 1.008 =
\qty{417.25}{g.mol^{-1}}$；两个 OH 合成一个 \ce{H2O}（$18.015$），占 \qty{4.32}{\percent}。
(b) \ce{K2O} $47.10/417.25 = \qty{11.3}{\percent}$；\ce{MgO} $3 \times 40.304/417.25 =
\qty{29.0}{\percent}$；\ce{Al2O3} $50.98/417.25 = \qty{12.2}{\percent}$；\ce{SiO2} $3 \times
60.083/417.25 = \qty{43.2}{\percent}$；合计 \qty{95.7}{\percent}。(c) \omterm{def:b2:field-lab-methods-2:epma}{电子探针}不能检测氢：
云母分析缺少其中的水，约 \qty{4}{\percent}，合计应接近
96。
\end{solution}

\begin{solution}{exo:b2:field-lab-methods-2:8}
(a) $\beta = \ln(0.1194/0.11905)/\ln(85.909261/87.905612) = 0.002936/(-0.02297) = -0.1278$。
(b) $0.71062 \times (86.908878/85.909261)^{-0.1278} = 0.70957$。(c) 测得的 86/88 比值低于
真实比值：重同位素被偏向，测得的 87/86 比值偏高。
\end{solution}

\begin{solution}{exo:b2:field-lab-methods-2:9}
(a) $0.1000 \times 15.271 \times 10^{-6}/83.9242 \times 1.3599 \times 10^{-3} = \qty{2.474e-11}{mol}$。
(b) 由 \eqref{eq:b2:field-lab-methods-2:id}，$N_x = 2.474 \times 10^{-11} \times (732.8 -
0.874)/(0.874 - 0.0568) = \qty{2.216e-8}{mol}$ 的 $^{86}\mathrm{Sr}$；锶：$2.216 \times
10^{-8}/0.0986 \times 87.62 = \qty{1.969e-5}{g}$，位于 \qty{0.0500}{g} 中，即 \qty{394}{ppm}（393.9）。
(c) 回收率为 \qty{99.95}{\percent}：该方法在比值取整的精度范围内是准确的。
\end{solution}

\begin{solution}{exo:b2:field-lab-methods-2:10}
(a) $R_m = \sqrt{732.8 \times 0.05680} = 6.45$，其中 $\mathcal M = 1.018$。(b) 在 0.1 处 $\mathcal M = 2.31$，
在 0.5 处为 1.13，在 1.664 处为 1.038，在 100 处为 1.16。(c) 在 $R_x$ 附近，混合物与样品几乎无异，
在 $R_s$ 附近则与稀释剂几乎无异，携带信息的差异淹没在 $R_m$ 的误差之中：$\mathcal M$ 在两端均发散。
两端之间的极小值十分平缓（约从 0.4 到 100，$\mathcal M < 1.2$），
因此分析者只需稀释剂的正确数量级，
粗略的浓度（例如由 XRF 提供）即可满足要求。
\end{solution}

\begin{solution}{exo:b2:field-lab-methods-2:11}
(a) $\sigma = \sqrt{0.62 \times 0.38/400} = 0.0243$：$62 \pm 2.4$。(b) $(65 - 62)/2.43 = 1.2$
个标准差：该岩石很可能为花岗闪长岩。(c) $\sigma = 0.015$ 需要 $0.62 \times
0.38/0.015^2 = 1047.1$，故取 1048 个长石点。
\end{solution}

\begin{solution}{exo:b2:field-lab-methods-2:12}
(a) 每个点独立地漏掉的概率为 $1 - p$：$(1-p)^n$。(b) $0.99^n \le 0.05$
得 $n \ge \ln 0.05/\ln 0.99 = 298.1$，故取 299 个点。(c) $n \ge \ln 0.05/\ln 0.9995 = 5990$：网格
计点法并不适用；稀有副矿物需通过扫描整个薄片来寻找。
\end{solution}

\begin{solution}{pb:b2:field-lab-methods-2:1}
\noindent\textbf{第一部分。}
\textbf{1.} 石英 $372/1200 = \qty{31.0}{\percent}$，碱性长石 29.0，斜长石 27.0，黑云母
9.0，白云母 3.0，不透明氧化物 1.0。
\textbf{2.} $\sqrt{0.31 \times 0.69/1200} = 0.0134$；白云母：$\sqrt{0.97/(0.03 \times 1200)} =
\qty{16.4}{\percent}$。
\textbf{3.} $M = 9 + 3 + 1 = \qty{13.0}{\percent}$；去掉白云母后 $M' = 13 - 3 = \qty{10.0}{\percent}$，
即浅色岩。
\textbf{4.} $Q + A + P = 1044$ 个点：Q $= 35.6$，A $= 33.3$，P $= 31.0$。
\textbf{5.} $324/672 = \qty{48.2}{\percent}$，Q 介于 20 与 60 之间：花岗岩；介于 35 与
65 之间为二长花岗岩。
\textbf{6.} $\sqrt{0.482 \times 0.518/672} = 0.0193$：35 与 65 两条边界分别距离
6.9 与 8.7 个标准差。该岩石名称可靠。
\textbf{7.} $0.31 \times 0.69/0.01^2 = 2139$ 个点。

\noindent\textbf{第二部分。}
\textbf{8.} $0.0051/0.8000 = \qty{0.64}{\percent}$。
\textbf{9.} $0.7949/8.7949 = 0.0904$，且 $72 \times 0.0904 = \qty{6.51}{\percent}$ \ce{SiO2} 位于
\omterm{def:b2:field-lab-methods-2:xrf}{熔融玻璃片}中：样品被稀释为原来的十一分之一。
\textbf{10.} $\bar I = 24.91$，$\bar c = 62.283$，$\sum(I - \bar I)^2 = 31.918$，$\sum(I - \bar I)(c -
\bar c) = 80.009$：$a = 2.5067$ 每 kcps，$b = 62.283 - 2.5067 \times 24.91 = -0.16$。
\textbf{11.} $2.5067 \times 28.63 - 0.16 = \qty{71.61}{\percent}$。
\textbf{12.} BCR-2 $-0.04$，AGV-2 $+0.05$，RGM-1 $-0.01$：均小于
方法的标准差（0.1 至 0.2），因此标准曲线在此范围内适用。
\textbf{13.} $71.61 + 27.15 + 0.64 = \qty{99.40}{\percent}$，介于 97 与 101 之间：无需重复测定。
\textbf{14.} $71.61 \times 100/(71.61 + 27.15) = \qty{72.51}{\percent}$。

\noindent\textbf{第三部分。}
\textbf{15.} $0.0986/0.8258 = 0.1194$：归一化值即天然比值。
\textbf{16.} $\beta = 0.007313/(-0.022972) = -0.3183$。
\textbf{17.} $0.71655 \times (86.908878/85.909261)^{-0.3183} = 0.71392$。
\textbf{18.} 测得的 86/88 比值与真实值相差 $0.11853/0.1194 - 1 = \qty{-0.73}{\percent}$，
跨越两个质量单位，约每个质量单位 \qty{0.36}{\percent}，偏向重同位素。
\textbf{19.} 稀释剂的 $^{86}\mathrm{Sr}/^{88}\mathrm{Sr}$ 为 0.7661：混合物的 86/88 比值不再是
天然的 0.1194，若将其固定为 0.1194，则会把稀释剂当作仪器偏差处理。

\noindent\textbf{第四部分。}
\textbf{20.} 各比值对 88 之和：$561.4 + 0.7661 + 0.1819 + 1 = 563.348$；$561.4/563.348 =
\qty{99.65}{\percent}$，且 $0.7661/563.348 = \qty{0.13599}{\percent}$。
\textbf{21.} $0.2500 \times 15.271 \times 10^{-6}/83.9242 = \qty{4.549e-8}{mol}$ 的锶，且
$\times 1.3599 \times 10^{-3} = \qty{6.186e-11}{mol}$ 的 $^{86}\mathrm{Sr}$。
\textbf{22.} $561.4/0.7661 = 732.8$，且 $0.0056/0.0986 = 0.05680$。
\textbf{23.} $N_x = 6.186 \times 10^{-11} \times (732.8 - 1.6812)/(1.6812 - 0.0568) =
\qty{2.784e-8}{mol}$。
\textbf{24.} $\mathcal M = 1.6812 \times 732.74/(731.12 \times 1.6244) = 1.04$；比值贡献
$1.04 \times \qty{0.01}{\percent}$，稀释剂贡献 $0.031/15.271 = \qty{0.20}{\percent}$；按平方和合成
得 \qty{0.20}{\percent}，主要由稀释剂的标定决定。
\textbf{25.} 锶：$2.784 \times 10^{-8}/0.0986 = \qty{2.824e-7}{mol}$，即 $2.824 \times
10^{-7} \times 87.62 = \qty{2.474e-5}{g}$，位于 \qty{0.1000}{g} 中。瓦尔登花岗岩含锶
\textbf{\qty[text-series-to-math=true]{247}{ppm}}，$\pm 0.5$。
\end{solution}
